% !TeX root = ../../Book5.tex
% 纲要标题：### C.4 旗簇、层与 D-模的几何方法
% 纲要位置：工作记录/计划纲要.md，第 4637 行。
% #### C.4.1 完全旗与射影直线
% #### C.4.2 线丛截面与最高权模
% #### C.4.3 微分算子模与中心关系
% ---
% 节标题由计划纲要.md 中附录 C 的第 4 条要点整理。

\section{旗簇、层与 D-模的几何方法}
\label{b5:sec:C04}

% 纲要范围：
% * 旗簇、层、D-模和几何化方法的进一步学习

最高权表示可以由几何对象上的截面构造。秩一时，旗簇只有两个仿射坐标图，线丛的截面与李代数微分算子都可以直接写出。这些计算有助于分清三种对象：几何空间、其上的层，以及作用于层截面的微分算子。

\subsection{完全旗与射影直线}
\label{b5:subsec:C04:01}

\begin{definition}\label{b5:def:complete-flag}
设 \(n\) 为非负整数，\(V\) 为 \(n\) 维复向量空间。满足
\[
0=F_0\subset F_1\subset\cdots\subset F_n=V,
\qquad\dim_{\mathbb C} F_i=i\quad(0\le i\le n),
\]
的子空间列，称为 \(V\) 的\term[wanquan-qi]{完全旗}[complete flag]。所有完全旗组成的代数簇称为\term[qi-cu]{旗簇}[flag variety]，记 \(\operatorname{Fl}(V)\)。
\end{definition}

其代数结构可以放在 \(\prod_{i=1}^{n-1}\operatorname{Gr}(i,V)\) 中，由相邻子空间包含的关联方程给出。\(\operatorname{GL}(V)\) 对旗传递作用：给每面旗选择适应它的基，换基矩阵就把标准旗送到所给旗。标准旗的稳定子群是可逆上三角矩阵群 \(B\)，故旗簇也是齐性空间 \(\operatorname{GL}_n/B\)。

\ProseParagraphBreak
当 \(n=2\) 时，一面旗只需选择一条直线，所以 \(\operatorname{Fl}(\mathbb C^2)=\mathbb P^1\)。上三角群有两个轨道：点 \([1:0]\) 与其补 \(\Bigl\{[z:1]\Bigr\}\cong\mathbb A^1\)。前者闭，后者的闭包为整个射影直线。这是 Schubert 分解的最小例子，两个轨道对应 \(S_2\) 的两个元素。这里所有闭包都光滑；高维旗簇的胞腔仍光滑，但胞腔的闭包可以奇异，不能把二者混称。

\subsection{线丛截面与最高权模}
\label{b5:subsec:C04:02}

射影直线用坐标 \(z=Y/X\) 与 \(w=X/Y\) 的两个仿射图覆盖。对整数 \(m\)，线丛 \(\mathcal O(m)\) 的正则截面可以写成一对多项式 \(a(z),b(w)\)，在交上满足
\[
b(w)=w^m a(1/w).
\]
这个相容条件是层的粘合要求，而不是任意选择两个无关的函数。

\begin{proposition}\label{b5:prop:projective-line-sections}
设 \(m\) 为整数，\(\mathbb P^1=\mathbb P^1_{\mathbb C}\)，\(\mathcal O(m)\) 为在坐标 \(z=Y/X\)、\(w=X/Y\) 的两图上按 \(b(w)=w^ma(1/w)\) 粘合截面的线丛，\(\mathbb C[X,Y]_m\) 为全次数 \(m\) 的齐次多项式空间。则
\[
H^0\Bigl(\mathbb P^1,\mathcal O(m)\Bigr)
\cong
\begin{cases}
\mathbb C[X,Y]_m,&m\ge0,\\
0,&m<0.
\end{cases}
\]
当 \(m\ge0\) 时，算子
\(e=X\partial_Y\)、\(f=Y\partial_X\)、\(h=X\partial_X-Y\partial_Y\)
给出 \(\mathfrak{sl}_2(\mathbb C)\) 作用，所得左模为最高权 \(m\) 的简单模 \(L_m\)。
\end{proposition}
\begin{proof}
若 \(a(z)=\sum_{j\ge0}a_jz^j\)，粘合条件给出
\(b(w)=\sum_j a_j w^{m-j}\)。两边都是多项式，恰当 \(a_j=0\) 对所有 \(j>m\) 成立。\(m<0\) 时全部系数为零；\(m\ge0\) 时对应齐次多项式
\(\sum_{j=0}^m a_jX^{m-j}Y^j\)。

三个算子保持全次数，直接计算得 \([h,e]=2e\)、\([h,f]=-2f\)、\([e,f]=h\)。\(X^m\) 是权 \(m\) 的最高向量，而反复作用 \(f\) 得到 \(X^{m-j}Y^j\) 的非零倍数，生成全空间。\(e\) 在除最高向量以外的基向量上系数非零，所以任一非零子模取最高向量后必含 \(X^m\)，进而为全模。因此它是 \(L_m\)。
\end{proof}

这里约定变量 \(X,Y\) 按自然二维表示变换。若把同一多项式空间解释成对偶空间上的函数，并令群通过逆变换拉回，便会出现对偶表示；比较几何文献时应先核对作用方向。

\ProseParagraphBreak
在 \(X\ne0\) 图中把齐次式写成 \(X^m a(z)\)，直接求导得到
\[
e=\partial_z,\qquad
h=m-2z\partial_z,\qquad
f=mz-z^2\partial_z.
\]
在另一图则变成
\[
e=mw-w^2\partial_w,\qquad
h=-m+2w\partial_w,\qquad
f=\partial_w.
\]
例如 \(e\) 的换图由 \(b(w)=w^ma(1/w)\) 推出
\(w^m a'(1/w)=mw\,b(w)-w^2b'(w)\)。三个算子的换图因此与线丛粘合相容，确实定义在几何截面上。

\subsection{微分算子模与中心关系}
\label{b5:subsec:C04:03}

对光滑复代数簇 \(X\)，将正则函数的乘法算子记为零阶算子；递归规定线性算子 \(P\) 阶数至多为 \(r\)，当且仅当对所有局部正则函数 \(a\)，交换子 \([P,a]\) 阶数至多为 \(r-1\)。它们按局部限制组成微分算子层 \(\mathcal D_X\)。带有 \(\mathcal D_X\) 左作用的层称为\term[D-mo]{\(D\)-模}[D-module]。该作用同时包含函数乘法和微分作用，二者必须满足 Leibniz 关系。

\begin{proposition}\label{b5:prop:affine-line-differential-operators}
设 \(z\) 为复仿射直线 \(\mathbb A^1_{\mathbb C}\) 的坐标，\(\partial=\mathrm{d}/\mathrm{d}z\) 为 \(\mathbb C[z]\) 上的求导算子。整体代数微分算子环 \(\Gamma(\mathbb A^1_{\mathbb C},\mathcal D_{\mathbb A^1})\) 是第一 Weyl 代数
\[
D=\mathbb C\langle z,\partial\rangle/(\partial z-z\partial-1).
\]
每个算子唯一写成有限和 \(\sum_{j=0}^r a_j(z)\partial^j\)，其中 \(r\ge0\) 为整数、\(a_j(z)\in\mathbb C[z]\)。
\end{proposition}
\begin{proof}
关系来自多项式求导的乘法法则。按算子阶数归纳：若 \(P\) 阶数至多 \(r\)，则 \([P,z]\) 阶数至多 \(r-1\)，归纳可写成 \(\sum_j b_j(z)\partial^j\)。令
\(Q=\sum_j b_j(z)\partial^{j+1}/(j+1)\)，利用
\([\partial^{j+1},z]=(j+1)\partial^j\)
得到 \([P-Q,z]=0\)。故 \(P-Q\) 与全部多项式乘法交换，必为乘以 \((P-Q)(1)\) 的算子。这证明所需表达存在。若一个这样的表达为零，依次作用于 \(1,z,z^2,\ldots\)，在已知较低系数为零之后，第 \(j\) 次给出 \(j!a_j(z)=0\)，故所有系数为零，证明唯一性。
\end{proof}

例如多项式模 \(\mathbb C[z]\) 可写成 \(D/D\partial\)。另一个循环模 \(D/Dz\) 由向量 \(\delta\) 生成，满足 \(z\delta=0\)，并有
\[
z\partial^j\delta=-j\partial^{j-1}\delta.
\]
它以 \(\delta,\partial\delta,\partial^2\delta,\ldots\) 为基，是集中在原点的微分算子模的代数模型。多项式模与这个模都在 \(\mathbb C\) 上无限维，但函数和微分算子在两者上的作用不同。

\ProseParagraphBreak
在线丛 \(\mathcal O(m)\) 上作用的微分算子组成扭曲的算子层 \(\mathcal D_{\mathcal O(m)}\)。在每个平凡化图上它与通常的 \(\mathcal D\) 同构；换平凡化时须用线丛过渡函数共轭算子，因此不同图上的写法并不只差坐标代换。上一小节的三个算子给出
\[
U(\mathfrak{sl}_2)\longrightarrow
\Gamma\Bigl(\mathbb P^1,\mathcal D_{\mathcal O(m)}\Bigr).
\]
直接展开可得
\(h^2+2h+4fe=m(m+2)\)，所以该映射通过中心商
\(U(\mathfrak{sl}_2)/\Bigl(\Omega-m(m+2)\Bigr)\)
分解。它把最高权参数、中心关系与同一个线丛的局部算子联系起来。

\ProseParagraphBreak
一般旗簇的局部化理论进一步研究何种中心商模能够由这样的算子层恢复；要得到范畴等价，还须规定中心参数的正则性、支配或反支配约定，以及所用 \(D\)-模的有限性条件。上面的中心分解计算本身没有证明这些等价。几何表示论还研究 \(D\)-模的好滤过、余切丛中的特征簇和 holonomic 条件；它们涉及微分方程的维数限制，与“全纯函数”这一概念不同。这里保留 holonomic 原文。关于这些方法与最高权表示的关系，参见\cite[Further reading, Geometric representation theory]{Kirillov2008LieGroups}；进入系统学习前，可先用本节的两个坐标图把层、扭曲和中心参数的区别算清楚。
